\subsection{Certifying the loss from an observable summary}
\label{sec:summarycertificate}

A finite controller may retain fills and a rounded filter while discarding
timestamps or finer reports. The lower bound in \cref{thm:computablepartial}
then concerns the declared finite observations; it need not be a lower bound
for an agent observing the richer history. The following construction closes
this comparison under explicit conditional-law bounds. It uses the robust
dynamic-programming foundations of \citet{iyengar2005,nilim2005}. The role
here is to certify a feasible execution controller against a richer
information class, without assuming that its summary is Markov or sufficient.

\begin{assumption}[Conditional action-law envelopes for an observable summary]
\label{ass:summaryenvelopes}
There are $K$ decisions under a supplied controlled law on standard Borel
history spaces. Let $H_k$ be the full observed history and let
$Z_k=\Psi_k(H_k)\in\mathcal Z_k$ be a finite summary, updated from the
preceding summary, action and new observations. The initial history is fixed.
Stage and terminal charges are bounded. The nonempty finite feasible action
set satisfies
\[
 A_k(H_k)=A_k(Z_k).
\]
In particular, the summary retains every unresolved reservation, instruction
status and session variable needed for feasibility. For every feasible
$(z,a)$, a nonempty compact set
$\mathcal U_k(z,a)\subset\mathbb R\times\Delta(\mathcal Z_{k+1})$
is supplied, and contains
\begin{equation}\label{eq:conditionalenvelope}
 \bigl(r_k(H,a),p_k(H,a)\bigr):=
 \bigl(\E[c_k\mid H,a],\ \mathcal L(Z_{k+1}\mid H,a)\bigr)
 \quad\text{whenever }\Psi_k(H)=z.
\end{equation}
Supplied terminal functions $l_K,u_K$ satisfy
$l_K(\Psi_K(H))\le\E[c_K\mid H]\le u_K(\Psi_K(H))$.
These inclusions hold for the conditional versions of the controlled model
on all admissible histories, or almost surely under every feasible policy,
including each feasible counterfactual action. Bounds established only under
one logging policy do not meet this assumption.
\end{assumption}

Write $p\cdot v=\sum_{z'}p(z')v(z')$. Define the finite recursions
\begin{align}\label{eq:summarysandwichdp}
 L_K&=l_K,\qquad U_K=u_K,\nonumber\\
 L_k(z)&=\min_{a\in A_k(z)}\inf_{(r,p)\in\mathcal U_k(z,a)}
                   \{r+p\cdot L_{k+1}\},\\
 U_k(z)&=\min_{a\in A_k(z)}\sup_{(r,p)\in\mathcal U_k(z,a)}
                   \{r+p\cdot U_{k+1}\}.\nonumber
\end{align}
Fix a minimizing action $\sigma_k(z)$ in the upper recursion, with a
deterministic tie rule. Compactness and finite action sets ensure attainment.
The recursions become numerical algorithms when their extrema can be
evaluated or enclosed, as in the two explicit constructions below.
Let $V_k^{\rm rich}(H)$ be the infimum of remaining expected cost over all
feasible policies using $H$ and subsequent full observations.

\begin{theorem}[A certificate against the richer history optimum]
\label{thm:summarysandwich}
Under \cref{ass:summaryenvelopes}, for $z=\Psi_k(H)$,
\begin{equation}\label{eq:summarysandwich}
 L_k(z)\le V_k^{\rm rich}(H)\le J_k^\sigma(H)\le U_k(z).
\end{equation}
In particular, with $z_0=\Psi_0(H_0)$,
\begin{equation}\label{eq:summaryregret}
 0\le J_\sigma-V_0^{\rm rich}(H_0)\le U_0(z_0)-L_0(z_0).
\end{equation}
For any other feasible summary controller $\rho$, including a randomized
finite-score signature controller, replace the minimum in the upper recursion
by the action average under $\rho_k(\cdot\mid z)$ and call the result
$U_k^\rho$. Then $J_\rho-V_0^{\rm rich}(H_0)\le U_0^\rho(z_0)-L_0(z_0)$.
Any independently certified upper evaluation $\overline E_\rho\ge J_\rho$
may replace $U_0^\rho(z_0)$.
\end{theorem}
\begin{proof}
At the terminal stage the asserted conditional bounds are the assumption.
Suppose the lower bound holds at $k+1$ for the conditional cost of every
feasible continuation. At history $H$, conditional on any chosen action $a$,
the tower property bounds that cost below by
$r_k(H,a)+p_k(H,a)\cdot L_{k+1}$. Membership in
\eqref{eq:conditionalenvelope} bounds it below by the corresponding infimum,
and then by $L_k(z)$. An average over randomized actions preserves the
inequality. Taking the infimum over all full-history continuations proves
$L_k(z)\le V_k^{\rm rich}(H)$, without restricting those continuations to
summary policies.

For the selected summary policy, the induction hypothesis and the tower
property instead give
\[
 J_k^\sigma(H)\le r_k(H,\sigma_k(z))+
       p_k(H,\sigma_k(z))\cdot U_{k+1}\le U_k(z).
\]
The selected action is feasible at every history in the summary fiber, so
this policy belongs to the class defining $V_k^{\rm rich}$. This proves the
middle inequality. For a fixed randomized controller, average the
actionwise suprema instead. Subtraction and an upper evaluation give the
remaining assertions.
\end{proof}

The stagewise envelope choices in \eqref{eq:summarysandwichdp} constitute
a rectangular bounding relaxation. A single physical model, or a fixed
unknown parameter, need not attain all those choices along one path.
The proof requires only conditional inclusion; it does not reset the hidden
state or assume that a physical adversary reselects the model at each step.
This distinction may make the certificate conservative.

\begin{corollary}[Exact value sufficiency and model transfer]
\label{cor:summarysufficiency}
If every envelope is the singleton
$\{(\widehat r_k(z,a),\widehat P_k(\cdot\mid z,a))\}$ and $l_K=u_K$,
then $L=U$ and the summary policy is optimal among all full-history policies.
More generally, suppose the envelope certificate is computed in a supplied
approximate model and
$|J_\pi-\widehat J_\pi|\le B_T$ uniformly over the common feasible
full-history policy class. Keep the summary update and the candidate $\rho$
fixed in both models. Then
\begin{equation}\label{eq:summarytransfer}
 J_\rho-\inf_{\pi\in\Pi_{\rm rich}}J_\pi
 \le 2B_T+\widehat{\overline E}_\rho-\widehat L_0(z_0).
\end{equation}
\end{corollary}
\begin{proof}
Singleton envelopes make both recursions the same Bellman recursion, so
\cref{thm:summarysandwich} forces equality. For transfer,
$J_\rho\le\widehat J_\rho+B_T$ and
$\inf_\pi J_\pi\ge\inf_\pi\widehat J_\pi-B_T$.
Apply the approximate-model certificate and subtract.
\end{proof}

\subsubsection*{A quantitative bound from conditional model discrepancies}

The envelopes themselves can retain state and action dependence. A simpler
uniform rate follows when they have the rectangular form
\begin{equation}\label{eq:summarytvenvelope}
 \begin{aligned}
 \mathcal U_k(z,a)=
 [\widehat r_k(z,a)-\alpha_k(z,a),\widehat r_k(z,a)+\alpha_k(z,a)]
 &\times\mathcal P_k(z,a),\\
 \mathcal P_k(z,a)=\{p\in\Delta(\mathcal Z_{k+1}):
              \|p-\widehat P_k(\cdot\mid z,a)\|_{\TV}
 &\le\epsilon_k(z,a)\},
 \end{aligned}
\end{equation}
where $\alpha_k\ge0$, $0\le\epsilon_k\le1$, and
$l_K=\widehat g-\delta_K$, $u_K=\widehat g+\delta_K$, $\delta_K\ge0$.
Let $\widehat V$ solve the nominal finite Bellman recursion with
$(\widehat r,\widehat P,\widehat g)$, and write
$\operatorname{span}(v)=\max v-\min v$.

\begin{proposition}[An explicit observation-reduction error bound]
\label{prop:summaryrate}
Define
\begin{align}\label{eq:summaryrateconstants}
 d_K&=\max_z\delta_K(z),\nonumber\\
 d_k&=\max_{z,a}\{\alpha_k(z,a)+
                   \epsilon_k(z,a)\operatorname{span}(\widehat V_{k+1})\},
 &B_k&=d_K+\sum_{j=k}^{K-1}d_j.
\end{align}
For the recursions \eqref{eq:summarysandwichdp},
\begin{equation}\label{eq:summaryrate}
 \widehat V_k-B_k\le L_k\le\widehat V_k\le U_k
                         \le\widehat V_k+B_k,\qquad
 U_0(z_0)-L_0(z_0)\le2B_0.
\end{equation}
The envelope recursions are finite linear programs. With rational data they
can be solved exactly; certified outward bounds on their extrema also give
a valid lower--upper certificate.
\end{proposition}
\begin{proof}
For probability vectors,
$|(p-\widehat P)\cdot v|\le\epsilon\operatorname{span}(v)$ under the
stated total-variation convention. Thus applying either envelope Bellman
operator to $\widehat V_{k+1}$ changes the nominal value by at most $d_k$.
Each operator is monotone and is one-Lipschitz in the sup norm, since its
transition vectors have unit mass. Backward induction from the terminal
error $d_K$ proves the outer bounds. The nominal pair belongs to every
envelope; a second induction gives $L_k\le\widehat V_k\le U_k$.
The absolute-value constraints defining a finite total-variation ball have
a linear-program representation. Lower bounds on the lower extrema and
upper bounds on the selected action's upper extrema propagate in the
required directions. An inexactly selected action is evaluated as a fixed
controller by \cref{thm:summarysandwich}.
\end{proof}

This bound uses nominal continuation-value spans, not a stability constant
for Bayes' formula. Computing the state-dependent recursions can retain
information discarded by the maxima in $d_k$. Neither refinement substitutes
for valid conditional error bounds.

\subsubsection*{A constructive witness for binary hidden liquidity}

Suppose the richer controlled model is hidden Markov with the binary state
of \cref{ass:finiteobservation}. A fine report is reduced to a symbol $o$,
and $M_{k,x,a,o}$ is its joint hidden-transition law after summing or
integrating the finer reports. The symbol retains the next ledger state
$F_k(x,a,o)$ and charge $c_k(x,a,o)$. The summary is $z=(x,\gamma)$ with
the fixed update \eqref{eq:beliefmemory}, driven by these coarse symbols.
Assume supplied intervals satisfying, for every admissible full history,
\begin{equation}\label{eq:posteriorfiber}
 \mathcal I_k(z)=[b_k^-(z),b_k^+(z)]\subset[0,1]
 ,\qquad \Pp(h_k=1\mid H_k)\in\mathcal I_k(\Psi_k(H_k)).
\end{equation}
This is a condition on the full posterior, not a claim that it equals
the controller's memory $\gamma$.

For each $b\in\mathcal I_k(z)$ define
\begin{align}\label{eq:binarysummaryenvelope}
 r_b(z,a)&=\sum_o p_b^a(o)c_k(x,a,o),\nonumber\\
 P_b(z'\mid z,a)&=\sum_o p_b^a(o)
 \mathbf 1\{z'=(F_k(x,a,o),Q_m\Phi_\gamma^a(o))\},\\
 \mathcal U_k(z,a)&=\{(r_b(z,a),P_b(\cdot\mid z,a)):
                                      b\in\mathcal I_k(z)\}.\nonumber
\end{align}
The update uses $\gamma$, while the physical probabilities use $b$.
The chosen zero-probability update convention is fixed for every symbol,
including a symbol impossible under $\gamma$ but possible under $b$.
Terminal bounds are the minimum and maximum of
$(1-b)c_K(x,0)+bc_K(x,1)$ on $\mathcal I_K(z)$.

\begin{proposition}[Endpoint evaluation of posterior envelopes]
\label{prop:binarysummary}
Under these hypotheses, \eqref{eq:binarysummaryenvelope} satisfies
\cref{ass:summaryenvelopes}. Each infimum and supremum in
\eqref{eq:summarysandwichdp} is attained at $b_k^-(z)$ or $b_k^+(z)$.
For supplied intervals, the backward recursions require
$O(K\max_k|\mathcal Z_k|\,|A|\,|O|)$ arithmetic operations.
If $|b-\gamma|\le\eta_k(z)$ throughout the interval, rectangular
envelopes of the form \eqref{eq:summarytvenvelope} are valid with
\begin{equation}\label{eq:binarysummaryradius}
 \alpha_k(z,a)=\eta_k(z)|\bar c_1-\bar c_0|,\qquad
 \epsilon_k(z,a)=\min\{1,\eta_k(z)\|P^1-P^0\|_{\TV}\},
\end{equation}
where $\bar c_h$ and $P^h$ are the conditional mean charge and next-summary
law given current hidden state $h$, $z$ and $a$. The nominal pair is the
mixture at $\gamma$. The analogous terminal radius is
$\eta_K(z)|c_K(x,1)-c_K(x,0)|$.
\end{proposition}
\begin{proof}
The controlled hidden-Markov property makes the conditional coarse outcome
law the mixture in \eqref{eq:finitebeliefkernel} at the full posterior $b$.
The memory update is already fixed by $z,a,o$. Therefore both the charge
mean and next-summary law in \eqref{eq:binarysummaryenvelope} are affine
in $b$, proving conditional inclusion and compactness. Pairing with any
fixed continuation vector remains affine, so the two endpoints suffice.
Each endpoint evaluation sums over at most $|O|$ symbols. Finally,
$r_b-r_\gamma=(b-\gamma)(\bar c_1-\bar c_0)$ and
$P_b-P_\gamma=(b-\gamma)(P^1-P^0)$, proving the radius formulas.
\end{proof}

The paired envelope preserves the dependence between charge and transition
errors through the same posterior. It can be tighter than their rectangular
enclosure. On a finite observation tree, intervals in
\eqref{eq:posteriorfiber} can be constructed by exact enumeration of all
positive-probability histories under all feasible actions, followed by
grouping at each summary. That preparatory enumeration can grow
exponentially; the stated arithmetic bound covers the subsequent backward
recursions, not the construction of the intervals. Continuous reports require
a separate analytic or validated numerical enclosure. The trivial interval
$[0,1]$ is always available but may yield a weak certificate.

In particular, $\eta_k(z)$ cannot generally be set to $1/(2m)$: repeated
rounding and discarded reports separate the full posterior from $\gamma$.
The one-step quantization radius in \cref{prop:partialmeshrate} belongs to
its resampling comparison and does not bound this different error.
Given valid posterior intervals, a controller evaluated by
\eqref{eq:truecontrollerevaluation} has the directly computable certificate
\[
 J_\rho-V_0^{\rm rich}(H_0)\le U_m(\rho)-L_0(z_0).
\]
Here the physical evaluation may integrate out the discarded reports because
the fixed controller does not use them. The new lower bound $L_0$ is what
makes the comparison valid against policies that do use those reports.
Section~\ref{sec:summaryexperiment} gives exact finite-tree checks and an
example where information reduction has a strictly positive cost.

\subsubsection*{Action comparisons can certify more than a cost interval}

An interval for the optimal cost can remain wide even when the same action
is optimal throughout a summary fiber. The next result uses the existing
conditional envelopes to certify decisions directly. The Bellman comparison
and upper/lower action screening are standard bounding principles; the
statement below keeps the richer-history information class and the exact
execution mask of \cref{ass:summaryenvelopes}.

For the lower and upper functions in \eqref{eq:summarysandwichdp}, define
\begin{align}\label{eq:summaryactionbounds}
 Q_k^-(z,a)&=\inf_{(r,p)\in\mathcal U_k(z,a)}
                     \{r+p\cdot L_{k+1}\},\nonumber\\
 Q_k^+(z,a)&=\sup_{(r,p)\in\mathcal U_k(z,a)}
                     \{r+p\cdot U_{k+1}\}.
\end{align}
If $a$ is the only feasible action, set $e_k(z,a)=0$. Otherwise set
\begin{equation}\label{eq:summaryactionexcess}
 e_k(z,a)=
 \left[Q_k^+(z,a)-\min_{a'\in A_k(z)\setminus\{a\}}Q_k^-(z,a')\right]_+.
\end{equation}
For any feasible summary controller $\rho$, use the finite recursion
\begin{align}\label{eq:summaryactionrecursion}
 R_K^\rho(z)&=0,\nonumber\\
 R_k^\rho(z)&=\sum_{a\in A_k(z)}\rho_k(a\mid z)
 \left\{e_k(z,a)+
       \sup_{(r,p)\in\mathcal U_k(z,a)}p\cdot R_{k+1}^\rho\right\}.
\end{align}
The supremum in this last display concerns continuation only; the immediate
cost is already accounted for through the action comparison.

\begin{theorem}[An action-comparison execution certificate]
\label{thm:summaryactioncertificate}
Under \cref{ass:summaryenvelopes}, every feasible summary controller satisfies
\begin{equation}\label{eq:summaryactioncertificate}
 0\le J_k^\rho(H)-V_k^{\rm rich}(H)
       \le R_k^\rho(\Psi_k(H)).
\end{equation}
Consequently an independently certified upper evaluation
$\overline E_\rho\ge J_\rho$ gives the combined bound
\begin{equation}\label{eq:summarycombinedcertificate}
 J_\rho-V_0^{\rm rich}(H_0)\le
 \min\{\overline E_\rho-L_0(z_0),\,R_0^\rho(z_0)\}.
\end{equation}
If every action used by $\rho$ has $e_k(z,a)=0$ at every stage and summary,
then $\rho$ is optimal among all richer-history policies. This conclusion
does not require singleton envelopes or equality of $L$ and $U$.
For the binary posterior envelopes, both action bounds and every step of
\eqref{eq:summaryactionrecursion} use only their two endpoints.
\end{theorem}
\begin{proof}
Let $Q_k^*(H,a)$ be the true expected charge plus the optimal richer-history
continuation after action $a$. The tower property, conditional inclusion
and \cref{thm:summarysandwich} give
\[
 Q_k^-(z,a)\le Q_k^*(H,a)\le Q_k^+(z,a).
\]
With at least two feasible actions, its Bellman advantage satisfies
\[
 Q_k^*(H,a)-V_k^{\rm rich}(H)
 =\left[Q_k^*(H,a)-\min_{a'\ne a}Q_k^*(H,a')\right]_+
 \le e_k(z,a).
\]
With one feasible action the advantage is zero. At $K$ the policy and
optimal conditional terminal costs agree. If the claimed regret bound
holds at $k+1$, condition on the action selected by $\rho$ and subtract the
true Bellman value. The resulting difference is the action advantage plus
the conditional expected continuation regret. Bound the first by $e_k$,
the second by $p_k(H,a)\cdot R_{k+1}^\rho$, and then use membership in the
conditional envelope. Averaging over $\rho_k$ proves
\eqref{eq:summaryactioncertificate} by backward induction.
The earlier cost-interval certificate gives the other endpoint in
\eqref{eq:summarycombinedcertificate}. If all used action excesses vanish,
the recursion gives $R^\rho=0$. Finally, the binary transition law is affine
in its posterior, so pairing it with any fixed continuation vector again
attains its extrema at the interval endpoints.
\end{proof}

The minimum of the two certificates is useful: neither bound is uniformly
sharper. A cost interval measures uncertainty in value, while
\eqref{eq:summaryactionexcess} compares the selected action with its
alternatives. The same bounds also give safe action screening. If
$Q_k^-(z,a)>\min_{a'} Q_k^+(z,a')$, action $a$ is strictly suboptimal at every
full history in that fiber and may be removed without changing the
richer-history optimum. The inequality is strict so that optimal ties are
not removed. This screening concerns optimization; it does not relax a
reservation, order-status or completion constraint.

These conclusions remain conditional on valid envelopes for all feasible
histories and actions. A small fitted action gap under one set of reports
does not establish the required uniform comparison. For model transfer,
compute the minimum in \eqref{eq:summarycombinedcertificate} in the
approximate model and add $2B_T$ as in \cref{cor:summarysufficiency}.
